\subsection{A Frobenius-orbit obstruction at inert prime powers}
\label{sec:frobenius-sidon}

\paragraph{Finding A24 (a necessary congruence for the A13 curve).}
Let $C/\mathbf F_2$ be the ordinary genus-two curve of A13,
\[
             y^2+y=x+x^{-1}+(x-1)^{-1},
\]
and let $X=C^{(1)}$, $J=\operatorname{Pic}^0(X)$ and $\Theta$
be its Raynaud theta support.  Let $n=\ell^r$, where $\ell$
is an odd prime and $r\ge1$, and assume the actual finite-level
condition
\[
                   \langle2,-1\rangle=(\mathbf Z/n)^\times.
                                                        \tag{SO.1}
\]
If $\Theta$ contains a point of exact order $n$, then $\ell$
is not inert in $\mathbf Q(\sqrt{-15})$.  Equivalently,
\[
        \ell\in\{3,5\}\quad\hbox{or}\quad
                    \ell\equiv1,2,4,8\pmod {15}.          \tag{SO.2}
\]
This is an unconditional geometric restriction.  In particular,
allowing higher powers satisfying (SO.1) does not restore an
inert underlying prime.  No assertion that every prime lifts to
all powers, and no general non-Wieferich hypothesis, is used.

The surviving arithmetic candidates are not exhausted by this
restriction.  Under premise M, the set of split primes satisfying
$\langle2,-1\rangle=(\mathbf Z/\ell)^\times$ has natural density
\[
                         \frac{72}{95}\mathcal A>0,        \tag{SO.3}
\]
where $\mathcal A$ is Artin's constant.  The smaller set of split
primes for which $2$ itself is primitive has density
$48\mathcal A/95$.  These conditional density statements supply
arithmetic candidates and do not assert theta incidence.

\begin{proof}[Proof of the unconditional restriction]
\emph{1. The finite-subgroup bound.}
For every nonzero $T\in J(\overline{\mathbf F}_2)$, A19 proves
\[
                   \#\bigl(\Theta\cap(\Theta+T)\bigr)\le2.
                                                        \tag{SO.4}
\]
Here and below these are numbers of geometric points.  If $H$ is
any finite subgroup of $J(\overline{\mathbf F}_2)$ and
$A\subseteq\Theta\cap H$, count the ordered pairs of distinct
points of $A$ by their nonzero difference.  Each difference has
at most two representations by (SO.4).  Thus
\[
                         |A|(|A|-1)\le2(|H|-1).           \tag{SO.5}
\]
In particular, for a prime-order point $z\in\Theta$ and its
Frobenius-generated $\mathbf F_\ell$-space $W$ of dimension $d$,
its Frobenius orbit $O$ satisfies
$|O|(|O|-1)\le2(\ell^d-1)$.  The same counting proof applies
to finite modules of prime-power exponent; it is not restricted
to vector spaces.

We will also use symmetry.  The characteristic-two locally exact
differential line bundle satisfies $B^{\otimes2}\simeq\omega_X$.
Riemann--Roch in degree one therefore gives
\[
 h^0(X,B\otimes L)=h^0(X,B\otimes L^{-1})
                        \qquad(L\in J),
\]
so $\Theta=-\Theta$.  This uses the same Raynaud line bundle as
in A03 and A11, not an arbitrary translate of the Abel curve.

\emph{2. The exact quadratic Frobenius relation.}
The equation and its Frobenius twist have the same model over
$\mathbf F_2$.  Let $\psi$ be $2$-power Frobenius and put
$\Phi=\psi^2$, the $4$-power Frobenius.  The checked A13 point
counts give the characteristic polynomial
\[
                    P_2(T)=T^4-T^2+4.
\]
Cayley--Hamilton on each odd Tate module gives the integral
identity
\[
                      \Phi^2-\Phi+4=0.                    \tag{SO.6}
\]
Thus the quadratic annihilator is established directly; it is
not inferred from the squared characteristic polynomial over
$\mathbf F_4$ without a semisimplicity argument.

Suppose henceforth that $\ell$ is inert.  Then
$Q(T)=T^2-T+4$ is irreducible modulo $\ell$, since its
discriminant is $-15$.  In particular $\ell\notin\{3,5\}$.
Let $\beta$ denote the class of $T$ in
$\mathbf F_{\ell^2}=\mathbf F_\ell[T]/(Q)$.  Its norm is
\[
                         \beta^{\ell+1}=4.                \tag{SO.7}
\]
On every nonzero vector of $J[\ell]$, the $\mathbf F_\ell[\Phi]$
span is a one-dimensional $\mathbf F_{\ell^2}$-space, hence a
two-dimensional $\mathbf F_\ell$-space, and the $\Phi$-orbit has
length
$s_1=\operatorname{ord}_{\mathbf F_{\ell^2}^{\times}}(\beta)$.

\emph{3. A uniform signed-orbit bound at the residue level.}
Write $h=(\ell-1)/2$.  Reduction of (SO.1) modulo $\ell$ gives
the signed-unit condition there.  In a cyclic group of even
order, adjoining $-1$ can enlarge a subgroup by at most two.
Consequently the order of $2$ is either $\ell-1$, or $h$ with
$h$ odd.  In either case
\[
                       \operatorname{ord}_\ell(4)=h.       \tag{SO.8}
\]
Let $H_1=\langle\beta,-1\rangle\subset\mathbf F_{\ell^2}^\times$.
Its norm image is $\langle4\rangle$, of order $h$, while the
norm kernel contains $-1$ and therefore has even order.
If this kernel had order two, then $|H_1|=2h=\ell-1$.
The unique subgroup of that order in the cyclic group
$\mathbf F_{\ell^2}^\times$ is $\mathbf F_\ell^\times$.
This would put $\beta$ in $\mathbf F_\ell$, contradicting
irreducibility.  The kernel consequently has order at least
four, and
\[
                          |H_1|\ge4h=2(\ell-1).            \tag{SO.9}
\]
This argument includes every inert eligible prime; it does not
require a large-prime cutoff or a list of exceptional cases.

\emph{4. Actual prime-power levels.}
Put $\mathcal O=\mathbf Z_\ell[T]/(Q)$, the ring of integers
of the unramified quadratic extension of $\mathbf Q_\ell$.
Identity (SO.6) makes $T_\ell J$ an $\mathcal O$-module.
The quadratic polynomial is irreducible over $\mathbf Q_\ell$,
so a nonzero element of $\mathcal O$ acts invertibly on
$T_\ell J\otimes\mathbf Q_\ell$.  The Tate module is therefore
torsion free over the discrete valuation ring $\mathcal O$,
and is free of rank two over it.  For
$R=\mathcal O/\ell^r\mathcal O$ this gives
\[
                              J[\ell^r]\simeq R^2.
\]
Let $z$ have exact additive order $\ell^r$.  In these coordinates
at least one component of $z$ is a unit of the local principal
ring $R$; otherwise its order would be at most $\ell^{r-1}$.
Thus $a\mapsto az$ is an injection $R\hookrightarrow R^2$.
The Frobenius-generated subgroup $H=Rz$ consequently has
exactly $\ell^{2r}$ points, and the $\Phi$-orbit of $z$ has
length $s_r=\operatorname{ord}_{R^\times}(\beta)$.

The reduction kernel $R^\times\to\mathbf F_{\ell^2}^\times$
is an $\ell$-group of exponent dividing $\ell^{r-1}$.
The latter assertion follows by repeated binomial expansion of
$1+\ell a$, with $\ell$ odd.  Therefore
\[
                           s_r=s_1\ell^b,
                         \qquad0\le b\le r-1.              \tag{SO.12}
\]
The norm $R^\times\to(\mathbf Z/\ell^r)^\times$ sends
$\beta$ to $4$.  The cyclic-unit argument used in (SO.8),
now applied to the actual condition (SO.1), gives
\[
                   \operatorname{ord}_{\ell^r}(4)
                         =h\ell^{r-1}.                    \tag{SO.13}
\]
This norm order divides $s_r$.  As $s_1$ is prime to $\ell$,
(SO.12)--(SO.13) force $b=r-1$.

Put $H_r=\langle\beta,-1\rangle\subset R^\times$.
If $s_1$ is odd, then $s_r=s_1\ell^{r-1}$ is odd, so adjoining
$-1$ doubles both orders.  If $s_1$ is even, then $-1$ already
belongs to both cyclic groups: the unique element of order two
in the units of this odd-residue-characteristic local ring is
$-1$.  To see uniqueness, $a^2=1$ implies
$(a-1)(a+1)=0$, and one of the factors is a unit because
$2$ is a unit; hence $a=1$ or $a=-1$.
It follows in both cases, with $k=\ell^{r-1}$, that
\[
                      |H_r|=k|H_1|\ge2(\ell-1)k.           \tag{SO.14}
\]
If $z\in\Theta$, Frobenius invariance and symmetry put its
entire signed orbit $H_rz$ in $\Theta\cap Rz$.  The injection
$R\hookrightarrow Rz$ makes its size exactly $|H_r|$.
Every inert prime here is at least seven, since $3,5$ ramify.
For $\ell\ge5$,
\[
\begin{split}
 |H_r|(|H_r|-1)
  &\ge k^2\,2(\ell-1)(2\ell-3)\\
  &>2\ell^2 k^2
   >2(\ell^{2r}-1).
\end{split}                                                \tag{SO.15}
\]
The middle inequality is
$2(\ell-1)(2\ell-3)-2\ell^2=2\ell^2-10\ell+6>0$.
The first follows from $k\ge1$ and
$(ak)(ak-1)\ge k^2a(a-1)$ for $a\ge1$.
This contradicts (SO.5), since $|Rz|=\ell^{2r}$.
This proves the stated
prime-power exclusion.  Quadratic reciprocity, or the two
quadratic characters of conductors $3$ and $5$, gives the split
classes $1,2,4,8$ modulo $15$, proving (SO.2).
\end{proof}

\paragraph{A finite consistency check, not a proof premise.}
Exact orders of $\beta$ in the three smallest inert residue fields
are
\[
\begin{array}{c|r|r|r}
 \ell&\operatorname{ord}(\beta)&s_1(s_1-1)&2(\ell^2-1)\\\hline
 7&24&552&96\\
 11&60&3540&240\\
 13&84&6972&336
\end{array}
\]
The accompanying \texttt{check\_A24\_orbits.py} verifies these
orders by quadratic-ring multiplication and an independent
$2$-by-$2$ multiplication-matrix implementation.  The uniform
norm-kernel proof above does not depend on this finite check.

\begin{proof}[Proof of the conditional residual densities]
Apply Lenstra's fixed finite Frobenius-condition theorem to
$W=\langle2,-1\rangle\subset\mathbf Q^\times$ and the
splitting condition in $F=\mathbf Q(\sqrt{-15})$.
The sufficient analytic family is
\[
 \zeta_{L_n}(s),\qquad
 L_n=\mathbf Q\bigl(\zeta_n,2^{1/n},(-1)^{1/n}\bigr),
                   \qquad n\ge1\text{ squarefree}.          \tag{SO.16}
\]
These Dedekind zeta functions are contained in premise M.
Lenstra's Section~2 permits an infinite finitely generated
subgroup with torsion, so it applies to this $W$ directly.
Theorem~3.1 and its Section~3.2 argument allow the fixed field
$F$ without requiring a separately enlarged GRH family
\cite[\S\S2--3, Theorem~3.1]{Lenstra1977}.

For odd squarefree $m$, write
$K_m=\mathbf Q(\zeta_m,2^{1/m})$.  As proved in A17,
\[
 [K_m:\mathbf Q]=m\varphi(m),\qquad
 K_m\cap\mathbf Q(\zeta_8)=\mathbf Q,
\]
and the maximal abelian subfield of $K_m$ is
$\mathbf Q(\zeta_m)$.  For completeness, Eisenstein at $2$
gives the degree, while the affine Galois group
$C_m\rtimes(\mathbf Z/m)^\times$ has commutator $C_m$:
the cyclotomic element $2$ conjugates a radical generator
$\tau$ to $\tau^2$.  The assertion about the abelian subfield
and the coprime-conductor intersection follow.

Taking odd roots of $-1$ adds nothing.  Taking a $2m$th root
of $-1$ adds $i$, and a $2m$th root of $2$ is generated by
$\sqrt2$ and $2^{1/m}$, using Bezout for $2,m$.  Thus
\[
 L_m=K_m,\qquad L_{2m}=\mathbf Q(\zeta_8)K_m,\qquad
 [L_m:\mathbf Q]=m\varphi(m),\quad
 [L_{2m}:\mathbf Q]=4m\varphi(m).                           \tag{SO.17}
\]
The maximal abelian subfields are respectively
$\mathbf Q(\zeta_m)$ and $\mathbf Q(\zeta_{8m})$.
The quadratic field $F$ has discriminant $-15$, is contained
in $\mathbf Q(\zeta_{15})$, and ramifies at both $3$ and $5$.
It follows that, for either $n=m$ or $n=2m$,
\[
 F\cap L_n=
 \begin{cases}
 F,&15\mid m,\\
 \mathbf Q,&15\nmid m.
 \end{cases}                                               \tag{SO.18}
\]
Indeed the ramification at a missing one of $3,5$ excludes
containment, while $15\mid m$ gives the containment directly.

For each finite collection of prime obstructions, ordinary
Chebotarev and inclusion-exclusion give the corresponding
finite sum of $\mu(n)/[FL_n:\mathbf Q]$.  Lenstra's theorem,
under (SO.16), identifies its limit with natural density.
The absolutely convergent sums can now be evaluated.  First
the total eligible density, without the splitting condition, is
\[
 D=\sum_{n\ {\rm squarefree}}\frac{\mu(n)}{[L_n:\mathbf Q]}
   =\frac34\prod_{q\ {\rm odd\ prime}}
                     \left(1-\frac1{q(q-1)}\right)
   =\frac32\mathcal A.
\]
Forcing $15\mid n$ changes the Euler factors at $3,5$ by
\[
 \frac{-1/6}{1-1/6}\cdot\frac{-1/20}{1-1/20}
                            =\frac1{95}.
\]
By (SO.18), the split density is therefore
\[
 \sum_n\frac{\mu(n)}{[FL_n:\mathbf Q]}
       =\frac12D+\frac12\frac{D}{95}
       =\frac{72}{95}\mathcal A.
\]
For $W=\langle2\rangle$ the even field is instead
$K_{2m}=K_m(\sqrt2)$ of degree $2m\varphi(m)$; the same
intersection argument applies.  The total density is then
$\mathcal A$, and its split part is $48\mathcal A/95$.
A sufficient GRH family for this smaller-set assertion is
$\zeta_{K_n}$ for squarefree $n$, as in A17.  Premise M
contains both stated sufficient families.
\end{proof}

\paragraph{Why the remaining candidates are not eliminated.}
The actual A13 witness $4159$ is split: $4159\equiv4\pmod{15}$.
Thus this same curve already prevents interpreting the
restriction as the absence of all eligible witnesses.
More generally, at a split prime away from the discriminant,
the quadratic Frobenius algebra is
$\mathbf F_\ell\times\mathbf F_\ell$.  For a point with
both eigencomponents nonzero, its generated space has dimension
two but its Frobenius orbit has at most $\ell-1$ points.
Such an orbit always satisfies the numerical inequality
$(\ell-1)(\ell-2)\le2(\ell^2-1)$, so this counting test
alone gives no contradiction.  An eigenline point can have a
shorter orbit; the A13 points over $\mathbf F_{128}$ have
$4$-power Frobenius orbit length dividing $7$.
The signed-unit condition controls the norm order; it does not
make theta invariant under multiplication by $2$.
Consequently A24 gives a necessary curve-dependent congruence
condition for prime and prime-power searches, while the
variable-prime theta-incidence gate of A18 remains open.
